\section{组员分工及检查点}
\label{sec:team}
\subsection{模块责任与协作方式}
三人分别承担自主模型、数据与适配实验、评价与分析，交付可运行模块及验证记录。罗开炜统筹集成，各模块由负责人实施、另一成员复核。

\begin{table}[!ht]
\caption{模块负责人、独立成果与交叉复核}
\label{tab:member-roles}\centering\small
\setlength{\tabcolsep}{4pt}\renewcommand{\arraystretch}{1.13}
\begin{tabularx}{\linewidth}{@{}P{1.35cm}YP{4.35cm}@{}}\toprule
\rowcolor{FLBlue!8}\textbf{成员} & \textbf{主责与独立技术成果} & \textbf{交叉复核}\\\midrule
\textbf{\color{FLBlue}罗开炜} & A/B/B-global、自主训练器、Hybrid 迁移与 SCST；训练恢复、推理接口及技术集成。 & 郝天乐复跑命令；张嘉荣核对模型图和对照。\\\addlinespace
\textbf{\color{FLTeal}郝天乐} & 数据审计、manifest、统一指标、GPU 排期；Qwen 数据处理、LoRA 训练与配对评测。 & 罗开炜审查监督掩码与适配范围；张嘉荣复算结果表。\\\addlinespace
\textbf{\color{FLAmber}张嘉荣} & 文献调研、Judge 接入与解析、人评协议、统计分析 Notebook；图表、报告与展示。 & 郝天乐复核统计输入与结果；罗开炜协助模型接入。\\\bottomrule
\end{tabularx}
\end{table}

郝天乐与张嘉荣独立盲评，罗开炜复核分歧。先以非测试样本校准标准，再锁定规则；正式评分隐藏候选来源及 Judge 结果。原始独立评分用于人评一致性与人工均分，协商裁定仅用于案例解释。抽样和统计协议与第五章一致。

\subsection{个人交付与可度量检查点}
各阶段按“产物、运行记录、通过条件、复核人”验收。表~\ref{tab:member-checks}明确个人成果，不以提交配置替代运行证据。
\begin{table}[!ht]
\caption{个人阶段交付与验收条件}
\label{tab:member-checks}\centering\small
\setlength{\tabcolsep}{4pt}\renewcommand{\arraystretch}{1.13}
\begin{tabularx}{\linewidth}{@{}P{2.15cm}YYY@{}}\toprule
\rowcolor{FLBlue!8}\textbf{检查点} & \textbf{\color{FLBlue}罗开炜} & \textbf{\color{FLTeal}郝天乐} & \textbf{\color{FLAmber}张嘉荣}\\\midrule
\textbf{开题}\newline 11 月 15 日 & 模型与接口表，明确输入、输出、迁移范围及选模依据。 & 数据来源与计数核验，图文映射、重复组及排除项审计办法。 & 文献与研究问题对应表，开题稿及课程任务覆盖核对。\\\addlinespace
\textbf{中期}\newline 12 月 6 日 & A/B 训练推理与 Hybrid 原型；郝天乐加载 A 权重复跑真实样本。 & manifest、哈希、GPU 试跑及四指标；预测 ID 与清单逐项对应。 & 独立复算指定指标并核对总表；同图案例、中期材料及人评草案。\\\addlinespace
\textbf{实验合并}\newline 12 月 25 日 & 迁移、SCST/等步数 CE 权重与验证记录；冻结测试版本。 & Qwen 配对结果和失败统计；罗开炜复跑适配器并核验监督范围。 & Judge、双人原始评分与统计 Notebook；郝天乐复算一致性及相关性。\\\addlinespace
\textbf{结题}\newline 1 月 17 日 & 核对代码、权重及配置；审查复现记录，统一提交。 & 全量预测、指标与资源记录；配合跨成员复算并解释差异。 & 报告及幻灯片或视频；独立运行演示，核对图表出处。\\\bottomrule
\end{tabularx}
\end{table}

\textbf{阶段牵头与交接。}开题、中期、结题分别由张嘉荣、郝天乐、罗开炜牵头，组长经教学云平台统一提交。每周核对进度，郝天乐登记 GPU 作业。交接包含代码版本、配置、数据哈希、命令、预测及失败统计；接收人核验 ID 全集与分母、复算指标。接口变动经交叉复跑后合并，阻塞时组长安排协助。

